\section{How to think about testing (from zero)}

Most people write tests by listing cases they happen to think of. That is layer-1
thinking and it gives false confidence. This section teaches the \emph{mental
models} that let you \emph{generate} tests systematically --- so you could build a
convincing suite for \emph{any} problem, and defend every test you wrote. The
concrete GoogleTest code follows in the next two sections; here we build the
\emph{judgment}.

\subsection{The core idea: a test is a trap for a specific bug}

A good test is not ``check it works.'' It is \textbf{a trap set for one particular
way the code could be wrong.} Before writing a test, finish this sentence:
\emph{``This test fails if the code \dots''}. If you can't name the bug, the test
isn't earning its place.

\begin{keybox}[The proof that weak tests lie]
Take a \emph{wrong} matching implementation --- the naive $\min(\text{buy},\text{sell})$
that ignores the company rule. Run it against ``obvious'' tests:
\begin{center}\footnotesize
\begin{tabular}{@{}lccc@{}}
\toprule
test & correct & naive-wrong & naive passes? \\
\midrule
partial (Buy 100 A, Sell 300 B) & $100$ & $100$ & \textcolor{Sell}{\textbf{yes}} \\
only buys & $0$ & $0$ & \textcolor{Sell}{\textbf{yes}} \\
simple cross & $100$ & $100$ & \textcolor{Sell}{\textbf{yes}} \\
self-only (Buy 100 A, Sell 100 A) & $0$ & $100$ & \textcolor{Buy}{no} \\
\textbf{dominant company} & $150$ & $950$ & \textcolor{Buy}{no} \\
\bottomrule
\end{tabular}
\end{center}
The broken code \textbf{passes three of five ``obvious'' tests.} Only the two tests
\emph{designed to discriminate} catch it. \emph{A suite that doesn't distinguish
right from a plausible wrong is theatre.}
\end{keybox}

\subsection{Mental model 1 --- read the tests off the formula}

Every term in a formula is a place the code can go wrong, so every term suggests a
test. Look at $M=\min(B,\ S,\ B+S-\max_c(b_c+s_c))$:

\begin{center}
\renewcommand{\arraystretch}{1.3}\footnotesize
\begin{tabular}{@{}p{3.2cm}p{5.3cm}p{5.0cm}@{}}
\toprule
\textbf{Term} & \textbf{Test that pins it} & \textbf{Fails if\dots} \\
\midrule
the $B$ (buy cap) & buys $<$ sells, cross-company & you cap by the wrong side \\
the $S$ (sell cap) & sells $<$ buys, cross-company & same, other side \\
the $-\max_c(\dots)$ term & \textbf{dominant company} & you used naive $\min$ \\
the different-company rule & self-only $\to 0$ & you match a firm with itself \\
\bottomrule
\end{tabular}
\end{center}

\begin{intu}
This is why the dominant-company test is \emph{the} test: it is the only one that
exercises the third, non-obvious term. If your suite has no test where one company
holds a big chunk of both sides, you have not tested the actual hard part.
\end{intu}

\subsection{Mental model 2 --- test at the boundaries, not the middle}

Bugs hide at edges, not in the comfortable interior. For every quantity, ask
``what are its extremes?'' and test each:

\begin{itemize}[leftmargin=1.5em,nosep]
  \item \textbf{Empty} --- no orders, unknown security/user/id. (Does it crash or
        return $0$/empty gracefully?)
  \item \textbf{One} --- a single order, a single company. (One company $\to$ always $0$.)
  \item \textbf{The threshold itself} --- \code{minQty} exactly equal to an order's
        qty. (\code{>} vs \code{>=} off-by-one.)
  \item \textbf{Huge} --- quantities near the type's max. (32-bit overflow.)
  \item \textbf{Duplicate} --- same id added twice. (Which policy fires?)
\end{itemize}

\begin{intu}
``Off-by-one'' and ``forgot the empty case'' are the two most common bugs in all
of software. Boundary tests are how you catch them on purpose instead of in
production. When you see \code{>=} in the code, you should reflexively write the
$=$ case.
\end{intu}

\subsection{Mental model 3 --- test relationships, not just values}

Fixed-value tests (``input $X$ gives $42$'') only cover cases you computed by hand.
\textbf{Metamorphic} tests check a \emph{relationship} that must hold for
\emph{every} input --- so they cover infinitely many cases you never enumerated.

\begin{center}
\renewcommand{\arraystretch}{1.3}\footnotesize
\begin{tabular}{@{}p{5.0cm}p{8.3cm}@{}}
\toprule
\textbf{Relationship (must always hold)} & \textbf{The bug it catches} \\
\midrule
add $X$ then cancel $X$ $\Rightarrow$ back to start & aggregate updated on add but not on cancel \\
insertion order doesn't change matching & hidden order-dependent state \\
cancel-for-user $\Rightarrow$ none of that user's orders remain & index desync / leftover ghosts \\
matching $\le \min(B,S)$ always & the cap logic is wrong \\
\bottomrule
\end{tabular}
\end{center}

\begin{intu}
The \emph{add-then-cancel round trip} is the single most valuable consistency test
here, because the \#1 bug in this design is ``I decrement the aggregate wrong on
cancel.'' If add-then-cancel returns you to the exact starting numbers, your two
sync helpers are mirror images --- which is the whole correctness argument.
\end{intu}

\subsection{Mental model 4 --- the oracle (test against a second brain)}

The most powerful idea in testing: compute the answer \textbf{two different ways}
and demand they agree. The second way --- the \textbf{oracle} --- must be written
\emph{differently} so it can't share a bug with the first.

\begin{center}
\begin{tikzpicture}[font=\footnotesize,>=Latex,node distance=8mm,every node/.style={align=center}]
  \node[draw,rounded corners,fill=Soft] (g) {random\\orders};
  \node[draw,rounded corners,fill=Buy!10,right=1.3cm of g] (i) {your cache\\(fast, maintained aggregate)};
  \node[draw,rounded corners,fill=Accent!12,below=8mm of i] (o) {oracle\\(slow, from raw orders)};
  \node[draw,rounded corners,fill=Gold!14,right=1.6cm of i] (c) {ASSERT\\equal};
  \draw[->](g)--(i); \draw[->](g)|-(o); \draw[->](i)--(c); \draw[->](o)-|(c);
\end{tikzpicture}
\end{center}

\begin{intu}
Why does this work? A bug in your fast code is unlikely to be \emph{the same} bug
in a slow, independently-written reference. So when they disagree on some random
input, one of them is wrong --- and you've found a case you'd never have dreamed up
by hand. Run it on thousands of random inputs (with a \textbf{fixed seed} so
failures reproduce) and you've effectively tested the space, not a handful of
points.
\end{intu}

\begin{pitfall}
The oracle must be \emph{obviously correct} and \emph{independent}. A subtle trap
here: a naive $O(n^2)$ greedy pairing is \emph{not} a correct oracle (greedy
under-counts, \S on math). The safe oracle is the per-company formula recomputed
from the \emph{raw order list} --- it re-derives the answer without reusing the
cache's maintained aggregate, so it catches maintenance bugs, which are the likely
ones.
\end{pitfall}

\subsection{Mental model 5 --- test the \emph{promise}, not just correctness}

The spec promised efficiency at a million orders. That is a claim your tests should
\emph{enforce}, because a correct-but-slow implementation (one that scans all orders
per query) passes every correctness test and still fails the real requirement.

\begin{intu}
A performance test is a correctness test for your \emph{design}. Insert a million
orders, time one \code{getMatchingSizeForSecurity}, assert it's under a few
milliseconds. A scanning implementation blows the budget; your aggregate
implementation sails through. The test is what proves the design paid off.
\end{intu}

\subsection{Putting the layers together (the pyramid)}

\begin{center}
\begin{tikzpicture}[font=\footnotesize]
  \node[draw,fill=Buy!10,minimum width=11.5cm,minimum height=0.8cm] (l1){\textbf{1. Examples} --- the spec's cases (ground truth)};
  \node[draw,fill=Accent!12,minimum width=9cm,minimum height=0.8cm,below=1mm of l1] (l2){\textbf{2. Edges} --- one trap per bug (esp.\ dominant company)};
  \node[draw,fill=Gold!14,minimum width=6.5cm,minimum height=0.8cm,below=1mm of l2] (l3){\textbf{3. Relationships} --- round trips, order-independence};
  \node[draw,fill=Sell!12,minimum width=4cm,minimum height=0.8cm,below=1mm of l3] (l4){\textbf{4. Oracle} + \textbf{perf}};
\end{tikzpicture}
\end{center}

\noindent Bottom-up: examples prove it works on known cases; edges catch the
specific bugs; relationships cover the unenumerated cases; the oracle covers the
\emph{space}; the perf test enforces the promise. A suite with all four reads as
\emph{engineering}, not homework --- which is exactly what a take-home grades.

\begin{say}
``I think of each test as a trap for one bug. I read the edge cases off the
formula --- the dominant-company case is the one that separates correct from naive.
I add relationship tests like add/cancel round-trips that catch aggregate-desync,
and a property test against an independent oracle on thousands of fixed-seed random
inputs so I'm testing the space, not points. Finally a million-order timing test
enforces the efficiency promise. I run it all under sanitizers so overflow and
iterator bugs fail loudly.''
\end{say}

\subsection{Blank-page drill --- generate the whole suite cold}

The real skill is producing the suite from \emph{nothing}. Close the document. On
a blank page, for each of the six methods, write every test you'd add \textbf{and
the one-line bug it traps}. Then uncover the checklist below and score yourself.

\begin{keybox}[Target: name the test \emph{and} the bug it catches]
\footnotesize
\renewcommand{\arraystretch}{1.25}
\begin{tabular}{@{}p{5.6cm}p{7.4cm}@{}}
\toprule
\textbf{Test you should recall} & \textbf{Bug it traps (say this too)} \\
\midrule
\multicolumn{2}{@{}l}{\textit{\color{Accent}Layer 1 --- examples}}\\
README example ($2700$) & baseline: the formula is plain wrong \\[2pt]
\multicolumn{2}{@{}l}{\textit{\color{Accent}Layer 2 --- edges (one trap each)}}\\
self-only $\to 0$ & a firm matches with itself \\
one cross-company match & counting own-company sells \\
\textbf{dominant company} $\to 150$ & using naive $\min(B,S)$ \\
only-one-side $\to 0$ & missing the empty-side guard \\
invalid orders ignored & no validation (empty id, qty 0, bad side) \\
duplicate id ignored & second add overwrites / corrupts \\
unknown lookup/cancel = no-op & missing graceful guard (throws) \\
\code{minQty} boundary ($=$) survives-check & \code{>} vs \code{>=} off-by-one \\
overflow (near \code{UINT\_MAX}) & 32-bit accumulator \\[2pt]
\multicolumn{2}{@{}l}{\textit{\color{Accent}Layer 3 --- relationships}}\\
add $\to$ cancel restores matching & aggregate not decremented on cancel \\
cancel-for-user leaves no trace & index desync / ghost ids \\
insertion order irrelevant & hidden order-dependent state \\[2pt]
\multicolumn{2}{@{}l}{\textit{\color{Accent}Layer 4 --- oracle + promise}}\\
random vs.\ oracle (fixed seed) & any matching bug, on cases you'd never dream up \\
random add/cancel vs.\ oracle & maintenance bug under churn \\
million-order timing & an $O(n)$ scan sneaking into the query \\
\bottomrule
\end{tabular}
\end{keybox}

\textbf{Score yourself:}
\begin{itemize}[leftmargin=1.5em,nosep]
  \item \textbf{$<8$ recalled} --- you're listing cases, not thinking in layers.
        Re-read the five mental models; the goal is to \emph{derive} these, not
        memorize them.
  \item \textbf{8--13} --- solid. Note \emph{which layer} you dropped (usually 3
        or 4 --- everyone remembers examples and edges).
  \item \textbf{14+} with the bug named for each --- you can defend a suite in the
        interview. Re-blank-page it tomorrow to make it permanent.
\end{itemize}

\begin{intu}
The trick to never blanking: don't memorize the 16 rows --- memorize the
\textbf{four layers} and the \textbf{five mental models}, then regenerate the rows
from them. ``Read tests off the formula'' alone gives you the whole edge layer;
``test relationships'' gives you layer 3; ``oracle + promise'' gives you layer 4.
The list is downstream of the thinking.
\end{intu}
